% Pixarch -- Baseline Report
% Compile:  tectonic docs/report.tex      (or: pdflatex docs/report.tex, twice)
\documentclass[11pt,a4paper]{article}

\usepackage[T1]{fontenc}
\usepackage[utf8]{inputenc}
\usepackage{lmodern}
\usepackage{amsmath}   % \text inside math mode
\usepackage[margin=2.2cm]{geometry}
\usepackage{booktabs}
\usepackage{longtable}
\usepackage{array}
% [table] pulls in colortbl, which provides \cellcolor (used by the caveat box).
\usepackage[table]{xcolor}
\usepackage{graphicx}
\usepackage{microtype}
\usepackage{enumitem}
\usepackage{fancyvrb}
% Long monospace filenames (IMG_20231205_123456.jpg) cannot hyphenate by
% default and run into the margin; [htt] permits breaking inside \texttt.
\usepackage[htt]{hyphenat}
\usepackage[hidelinks,colorlinks=true,linkcolor=linknavy,urlcolor=linknavy,
            citecolor=linknavy]{hyperref}

\definecolor{linknavy}{HTML}{1F4E79}
\definecolor{rulegrey}{HTML}{CCCCCC}
\definecolor{boxbg}{HTML}{F4F6F8}

% Callout box for caveats, built from primitives so no extra package is needed.
\newenvironment{caveat}
  {\par\vspace{0.6em}\noindent
   \begin{tabular}{@{}p{0.4em}@{\hspace{0.8em}}p{\dimexpr\linewidth-2.2em\relax}@{}}
   \cellcolor{linknavy}\rule{0pt}{1em} &}
  {\end{tabular}\par\vspace{0.6em}}

\newcommand{\code}[1]{\texttt{\small #1}}
\newcommand{\best}[1]{\textbf{#1}}

\setlength{\tabcolsep}{5pt}
\renewcommand{\arraystretch}{1.18}

\title{\vspace{-1.4cm}\textbf{Pixarch --- Baseline Report}\\[0.35em]
       \large On-device agentic retrieval over personal data}
\author{}
\date{31 October 2026}

\begin{document}
\maketitle
\vspace{-0.9cm}

\noindent\rule{\linewidth}{0.4pt}
\vspace{0.4em}

\noindent
This report covers the \textbf{baseline}: the non-agentic, single-shot system
that every later contribution is measured against. All numbers come from
\code{docs/results/baseline\_2026-10-31.md}, which \code{scripts/evaluate.py}
regenerates from scratch.

\begin{caveat}
\textbf{Status: indicative, not publishable.} Relevance is single-judge and part
of the corpus is generated. The numbers are good enough to catch regressions and
to size the on-device problem; they are not good enough to support a claim.
\end{caveat}

\noindent\rule{\linewidth}{0.4pt}

\tableofcontents
\vspace{0.6em}
\noindent\rule{\linewidth}{0.4pt}

%======================================================================
\section{Labeled Dataset Creation}

No dataset exists for mixed personal-archive search, so we built one. The corpus
grew from 127 items to \textbf{629 manifest rows / 627 indexed} across
\textbf{31 subtypes} --- 14 document, 17 scene.

\code{data/raw/} is gitignored while \code{data/manifest.csv} is committed,
carrying \code{source\_url}, \code{license} and \code{attribution} per row, so the
corpus is reproducible \emph{by reference} rather than by redistributing images.
\code{scripts/fetch\_dataset.py} rebuilds it.

\subsection{Sources}

\begin{longtable}{@{}p{1.8cm}p{6.6cm}p{4.4cm}p{2.6cm}@{}}
\toprule
\textbf{Bucket} & \textbf{Subtypes} & \textbf{Source} & \textbf{Licence} \\
\midrule
\endhead
scene & beach, mountain, trees, vehicle, party, wedding, college, food, travel,
        fitness, mall, planes, outdoor\_games, clothes, art, transit\_hub
      & SUN397 (via HuggingFace) & research use \\
scene & xray & keremberke/chest-xray-classification & CC-BY-4.0 \\
document & receipt & ICDAR-2019 SROIE & CC-BY-4.0 \\
document & bill & katanaml-org/invoices-donut-data-v1 & CC-BY-4.0 \\
document & general\_doc & FUNSD & research use \\
document & handwritten & IAM Handwriting Database & research use \\
document & book & Wikimedia Commons title pages & PD / CC \\
document & aadhaar, pan, dl, passport, credit\_card, debit\_card, passbook,
           bank\_statement, report\_card
         & \textbf{synthetic} (generated) & CC0 \\
\bottomrule
\end{longtable}

\subsection{Why nine classes are synthetic}

Wikimedia Commons holds \textbf{nine} usable Indian identity specimens in total
--- verified by search, not assumed. Beyond those, every image of an Aadhaar
card, PAN card, passport, licence, bank passbook or payment card on the open web
is a photograph of a \textbf{real person's} document, carrying their name,
address, photograph and the number itself.

This is not merely a privacy preference; it is a property of the repository.
\code{manifest.csv} is committed and stores \code{source\_url}, and
\code{fetch\_dataset.py} reconstructs the corpus from it. Fetching such images
would publish a permanent, public index of real people's identity documents and
cause everyone who clones the repository to re-download them. India's Aadhaar
Act 2016 additionally restricts the sharing of identity information. Building the
evaluation set would be worse still: OCR ground truth would mean transcribing
real Aadhaar numbers into a committed file.

Generation avoids all of this and produces \emph{better} data: exact ground
truth, balanced classes (21 PAN instead of 1), and controllable difficulty.

\paragraph{Safety properties of the generated values.}
\begin{itemize}[leftmargin=1.3em,itemsep=0.15em]
  \item Aadhaar numbers begin with \textbf{0 or 1} --- UIDAI never issues those,
        so no generated number can collide with a real one.
  \item Card numbers are constructed to \textbf{fail the Luhn checksum}, so none
        can be a valid payment card.
  \item All person, bank and school names are fictional.
\end{itemize}

There is deliberately \textbf{no ``SPECIMEN'' watermark}. A token appearing in
every synthetic document and no real one is a label leak straight into BM25 ---
the same defect as the filename leak that \code{schema.sql} documents removing
(worth $+0.125$ p@1, measured, then deleted).

The nine real Commons specimens are kept at their original indices and excluded
from the synthetic pool, so they remain a held-out realism check.

\subsection{Handwriting}

Every page-level handwriting dataset on HuggingFace is gated. Only line-level
sets are reachable, and a raw IAM line is ${\sim}2467\times128$ --- a 19:1 strip
that is neither a note a user would photograph nor anything CLIP can route
sensibly. Pages are therefore \textbf{composed from real IAM handwriting lines}:
the pixels are real, only the layout is ours, and the ground truth is the
concatenated line text.

\subsection{Ground truth and judgements}

\begin{center}
\begin{tabular}{@{}llrr@{}}
\toprule
\textbf{Artifact} & \textbf{Contents} & \textbf{Before} & \textbf{After} \\
\midrule
\code{ocr\_groundtruth.jsonl} & character-exact + upstream transcriptions & 20 & \best{260} \\
\code{queries.jsonl}          & $\geq 3$ per subtype, all 31 covered      & 20 & \best{98}  \\
\code{qrels.csv}              & judgements, expanded by subtype           & 185 & \best{2{,}183} \\
\bottomrule
\end{tabular}
\end{center}

\code{scripts/make\_qrels.py} now fails loudly if any corpus subtype has no
query --- the check that would have caught half the corpus becoming unevaluable
when the subtype count went $15 \rightarrow 31$.

%======================================================================
\section{Models used for PC simulation}

\begin{center}
\begin{tabular}{@{}p{5.2cm}p{4.6cm}p{5.6cm}@{}}
\toprule
\textbf{Role} & \textbf{Model} & \textbf{Footprint} \\
\midrule
Scene embedding + bucket routing & CLIP \code{ViT-B-32} / \code{laion2b\_s34b\_b79k}
  & \textbf{151.3\,M params} --- 605\,MB fp32, 303\,MB fp16, 151\,MB int8 \\
OCR & Tesseract 5, \code{eng+hin} & system package \\
Lexical index & SQLite FTS5 + BM25 & in-index, no model \\
\bottomrule
\end{tabular}
\end{center}

One CLIP instance serves \textbf{two} jobs --- zero-shot bucket classification at
ingestion and scene retrieval at query time --- which halves what the on-device
build must budget for.

Tesseract is the engine behind every number in this report. EasyOCR remains
implemented behind the same protocol as a comparator but is not used here.

%======================================================================
\section{Finalized models for the edge device}

\begin{center}
\begin{tabular}{@{}lll@{}}
\toprule
\textbf{Stage} & \textbf{Desktop} & \textbf{Android} \\
\midrule
OCR       & Tesseract \code{eng+hin} & \textbf{ML Kit Text Recognition v2} (Latin + Devanagari) \\
Embedding & CLIP ViT-B-32 fp32       & same weights, fp16 / int8 \\
Index     & SQLite FTS5              & SQLite FTS5 (ships with Android) \\
\bottomrule
\end{tabular}
\end{center}

All engines sit behind \code{OCREngine.extract(path) -> list[TextBlock]}, so the
Android build swaps OCR by changing one configuration key. FTS5 was chosen partly
because it is present on Android unchanged.

\paragraph{Budget implied by the measurements.}
CLIP at int8 is ${\sim}151$\,MB resident, and the index projects to 85\,MB for a
10{,}000-item gallery (\S\ref{sec:storage}). Ingestion is the real constraint:
1.56 hours for 10{,}000 items on a desktop GPU, which is why the architecture
assumes heavy processing runs while the phone is idle and charging.

%======================================================================
\section{Baseline Architecture}

\begin{center}
\begin{BVerbatim}
image --> ingest (hash, EXIF) --> (1) CLIP zero-shot classify
                                        |
                        +---------------+---------------+
                   document                          scene
                        |                               |
              OCR -> doctype rules            CLIP image embedding
                        |                               |
            FTS5 inverted index                   vector store
                        |                               |
                      BM25                           cosine
\end{BVerbatim}
\end{center}

At query time a \textbf{rule-first router} picks exactly one branch; CLIP
prototype similarity is consulted only when the keyword rules abstain. Everything
lands in a single \code{index.db}.

\paragraph{Deliberately absent.}
Each of the following is a later contribution that must not be pre-empted by the
baseline it will be measured against: no score fusion, no agentic loop, no retry
or clarification, no face branch.

\paragraph{Two choices the measurements justify.}
\begin{itemize}[leftmargin=1.3em,itemsep=0.2em]
  \item \code{doc\_type} is its own FTS column \textbf{weighted $5\times$} against
        body text. Raw BM25 over OCR text is brittle when a scan is noisy and the
        one distinguishing phrase comes back mangled.
  \item \textbf{No filename column.} Indexing generated names leaked the class
        label into the ranking --- worth $+0.125$ p@1, measured, then removed.
\end{itemize}

%======================================================================
\section{Results}

Corpus 627 items $\cdot$ 98 queries $\cdot$ retrieval depth 107 $\cdot$ CLIP on
CUDA $\cdot$ Tesseract.

\subsection{Retrieval}

\begin{center}
\begin{tabular}{@{}lrrrrrr@{}}
\toprule
\textbf{System} & \textbf{p@1} & \textbf{p@5} & \textbf{MRR@10} & \textbf{nDCG@10}
  & \textbf{R-prec} & \textbf{recall@10} \\
\midrule
\textbf{routed (baseline)} & \best{0.773} & \best{0.781} & \best{0.803} & \best{0.773} & \best{0.693} & 0.369 \\
clip\_only  & 0.629 & 0.596 & 0.644 & 0.580 & 0.478 & 0.278 \\
bm25\_only  & 0.361 & 0.355 & 0.378 & 0.349 & 0.312 & 0.157 \\
filename    & 0.299 & 0.299 & 0.299 & 0.299 & 0.298 & 0.140 \\
\bottomrule
\end{tabular}
\end{center}

Routing beats both single-model comparators overall: \textbf{+0.144 p@1} over
CLIP-only and \textbf{+0.412} over BM25-only.

\paragraph{Per bucket.}
Macro-averaged \emph{within} each bucket over its queries, then macro-averaged
\emph{across} buckets so each counts once regardless of query count.

\begin{center}
\begin{tabular}{@{}lrrr@{}}
\toprule
\textbf{System} & \textbf{document p@1} & \textbf{scene p@1} & \textbf{macro p@1} \\
\midrule
routed     & 0.739 & 0.804 & \best{0.772} \\
clip\_only & 0.370 & \best{0.863} & 0.616 \\
bm25\_only & \best{0.739} & 0.020 & 0.379 \\
\bottomrule
\end{tabular}
\end{center}

The honest reading: routing's value is that it is \emph{never catastrophic}. Each
single-model system collapses on the bucket it was not built for --- BM25 scores
0.020 on scenes, CLIP 0.370 on documents --- while routing tracks the better of
the two everywhere. On scenes alone CLIP-only is \emph{better} than routed
(0.863 vs 0.804), because every routing error is a pure loss.

\subsection{Recall with ceilings}

Recall@$k$ cannot exceed $k/|\text{relevant}|$. With ${\sim}20$ items per subtype
recall@10 caps at 0.48, so the raw figure badly understates performance.

\begin{center}
\begin{tabular}{@{}lrrr@{}}
\toprule
\textbf{System} & \textbf{recall@10} & \textbf{ceiling@10} & \textbf{\% of achievable} \\
\midrule
routed     & 0.369 & 0.480 & \best{77.0\%} \\
clip\_only & 0.278 & 0.480 & 57.9\% \\
bm25\_only & 0.157 & 0.480 & 32.7\% \\
filename   & 0.140 & 0.480 & 29.2\% \\
\bottomrule
\end{tabular}
\end{center}

R-precision --- self-normalising, with no ceiling artefact --- tells the same
story: \textbf{0.693} for routed against 0.478 for CLIP-only.

\subsection{Module benchmarks}
\label{sec:modules}

\paragraph{Bucket classification (CLIP zero-shot, ingestion).}
618/627 $=$ \textbf{98.6\%}. Eight scene $\rightarrow$ document, one
document $\rightarrow$ scene.

\paragraph{Query routing (rules, then CLIP).}
\textbf{95.9\%} --- 4 of 98 misrouted. Three are the x-ray queries: the router
has no radiography vocabulary, so CLIP decides, and a chest film reads as ``a
scanned document''. The fourth is the intentionally unanswerable \code{q020}.

\paragraph{Document subtype tagging (\code{doctype.py}) --- 70.7\% over 287
documents.} This module had never been measured.

\begin{center}
\begin{tabular}{@{}lrr@{}}
\toprule
\textbf{Origin} & \textbf{n} & \textbf{Accuracy} \\
\midrule
synthetic & 180 & \best{100.0\%} \\
fetched   & 107 & \best{21.5\%} \\
\bottomrule
\end{tabular}
\end{center}

\begin{center}
\begin{tabular}{@{}lrp{8.6cm}@{}}
\toprule
\textbf{Misclassification} & \textbf{n} & \textbf{Cause} \\
\midrule
bill $\rightarrow$ receipt & 20 & invoices carry ``total''/``invoice no''; no rule distinguishes them from receipts \\
book $\rightarrow$ unknown & 20 & \code{book-isbn} needs isbn/chapter/copyright; title pages have none \\
handwritten $\rightarrow$ unknown & 20 & no rule exists, and OCR on handwriting is weak \\
general\_doc $\rightarrow$ unknown/receipt & 20 & generic forms have no distinguishing vocabulary \\
aadhaar/pan/passport $\rightarrow$ unknown & 4 & OCR failed to recover the key phrase \\
\bottomrule
\end{tabular}
\end{center}

The $100\%$ / $21.5\%$ split is the headline: the rules are effectively
\emph{perfect on rendered text and poor on photographs}, so the aggregate
$70.7\%$ is an artefact of the synthetic:fetched ratio rather than a property of
the module. \code{receipt} precision of 0.391 --- recall 1.000, but 20 bills and
8 forms swept in --- is the clearest single defect.

\paragraph{OCR accuracy (Tesseract, 258 items with ground truth).}

\begin{center}
\begin{tabular}{@{}lrrr@{}}
\toprule
\textbf{Set} & \textbf{n} & \textbf{CER} & \textbf{WER} \\
\midrule
all                  & 258 & 0.471 & 0.596 \\
excluding \code{bill}& 238 & \best{0.420} & \best{0.527} \\
synthetic            & 180 & 0.418 & 0.507 \\
fetched              &  78 & 0.592 & 0.801 \\
\bottomrule
\end{tabular}
\end{center}

\begin{caveat}
\code{bill} is \textbf{not comparable}. The upstream invoice ground truth is a
set of extracted \emph{field values}, not a page transcription, so OCR that
correctly reads the whole page is scored against a much shorter reference and its
CER exceeds 1.0. This is a ground-truth mismatch, not an OCR failure.
\end{caveat}

The synthetic-versus-fetched gap (0.418 vs 0.592 CER) quantifies the domain gap
the generated corpus carries: rendered text is roughly 30\% easier than
photographs and scans. The best subtype is \code{passbook} (0.187); the worst
comparable one is \code{bank\_statement} (0.655).

\subsection{Query latency}

\begin{center}
\begin{tabular}{@{}llrrrrr@{}}
\toprule
\textbf{System} & \textbf{Scope} & \textbf{n} & \textbf{mean} & \textbf{p50}
  & \textbf{p95} & \textbf{max} \\
\midrule
routed     & all      & 98 & 8.0\,ms  & 6.8  & 13.5 & 63.3 \\
routed     & document & 49 & \best{4.1\,ms} & 1.6 & 8.5 & 46.9 \\
routed     & scene    & 49 & \best{11.9\,ms} & 12.4 & 15.2 & 63.3 \\
bm25\_only & ---      & 98 & 0.6\,ms  & 0.4  & 1.5  & 2.2 \\
\bottomrule
\end{tabular}
\end{center}

The branches are bimodal by an order of magnitude: BM25 is a SQL query
(0.6\,ms), while the scene branch pays a CLIP text-encoder forward pass
(11.9\,ms). Both sit far below interactive thresholds --- \textbf{query time is
not the on-device constraint}.

\subsection{Ingestion throughput}

\begin{center}
\begin{tabular}{@{}lrrr@{}}
\toprule
\textbf{Stage} & \textbf{Items} & \textbf{ms/image} & \textbf{images/sec} \\
\midrule
ingest        & 627 & 29.1  & 34.41 \\
clip\_encode  & 627 & 43.2  & 23.14 \\
classify      & 627 & 26.9  & 37.22 \\
embed\_store  & 627 & 0.2   & 4209.58 \\
\textbf{ocr}  & 393 & \best{738.0} & \best{1.35} \\
\midrule
\textbf{end-to-end} & 627 & \best{562.0} & \best{1.78} \\
\bottomrule
\end{tabular}
\end{center}

Stages run over different denominators: \code{ocr} touches 393 items (287
documents plus 106 low-confidence items routed to both branches) while the rest
touch all 627. A single ms/image averaged across stages would be meaningless.

\begin{center}
\begin{tabular}{@{}lrr@{}}
\toprule
\textbf{Bucket} & \textbf{ms/image} & \textbf{images/sec} \\
\midrule
scene    & 99.4   & 10.06 \\
document & 1110.0 & 0.90 \\
\bottomrule
\end{tabular}
\end{center}

\textbf{OCR is 82\% of total indexing time} and makes documents ${\sim}11\times$
more expensive than scenes, so mixed-gallery throughput depends entirely on the
document:scene mix.

Per-item OCR is long-tailed --- mean 691\,ms, p50 478\,ms, \textbf{p95
1567\,ms}, \textbf{max 3048\,ms} --- so the mean understates the worst case by
roughly $4\times$.

\begin{center}
\fcolorbox{rulegrey}{boxbg}{%
\begin{minipage}{0.92\linewidth}
\vspace{0.3em}
\centering \textbf{Projected: 1.56 hours to index a 10{,}000-item gallery} at
this corpus's 46\% document ratio (linear projection, desktop GPU).
\vspace{0.3em}
\end{minipage}}
\end{center}

\subsection{Storage and memory}
\label{sec:storage}

\code{index.db} is \textbf{5.3\,MB} for 627 items.

\begin{center}
\begin{tabular}{@{}lrr@{}}
\toprule
\textbf{Component} & \textbf{MB} & \textbf{Share} \\
\midrule
embeddings                 & 2.58 & 48.3\% \\
ocr\_blocks (+ autoindex)  & 1.54 & 28.9\% \\
doc\_fields                & 0.39 & 7.4\% \\
doc\_fts (inverted index)  & 0.36 & 6.7\% \\
items                      & 0.25 & 4.7\% \\
\bottomrule
\end{tabular}
\hspace{0.8em}
\begin{tabular}{@{}lr@{}}
\toprule
\textbf{Measure} & \textbf{Bytes} \\
\midrule
per item, overall                      & 8{,}525 \\
per item, shared (row + CLIP vector)   & 5{,}840 \\
extra per document (OCR + FTS)         & 5{,}866 \\
\bottomrule
\end{tabular}
\end{center}

\paragraph{Projected index for 10{,}000 items:}
85\,MB at this corpus's ratio; 58\,MB all scene; 117\,MB all document. With CLIP
at int8 ($\sim$151\,MB resident), the whole system fits comfortably inside a
mid-range phone's budget --- the storage claim holds.

%======================================================================
\section{Repository}

\begin{Verbatim}[fontsize=\small,frame=single,framesep=6pt,rulecolor=\color{rulegrey}]
pip install -r requirements.txt
python scripts/fetch_dataset.py --per-class 20 --per-doc 20   # rebuild corpus
python scripts/make_qrels.py                                  # queries -> qrels
python scripts/build_index.py                                 # ingest + OCR + CLIP
python scripts/evaluate.py                                    # -> docs/results/
python app/demo.py                                            # Gradio UI
pytest                                                        # 94 tests
\end{Verbatim}

\begin{center}
\begin{tabular}{@{}p{7.1cm}p{7.9cm}@{}}
\toprule
\textbf{Path} & \textbf{Contents} \\
\midrule
\code{src/pixarch/\{ingest,\allowbreak classify,\allowbreak ocr,\allowbreak doctype\}.py} & ingestion pipeline \\
\code{src/pixarch/synthetic.py}  & synthetic specimen generation \\
\code{src/pixarch/store/}        & SQLite schema, FTS5, vector store \\
\code{src/pixarch/query/}        & router, BM25, CLIP search \\
\code{src/pixarch/eval/}         & metrics, module benchmarks, resource accounting \\
\code{configs/baseline.yaml}     & every tunable \\
\code{docs/design-decisions.md}  & why each choice, and the traps hit \\
\bottomrule
\end{tabular}
\end{center}

%======================================================================
\section{Known limitations}

\begin{enumerate}[leftmargin=1.5em,itemsep=0.35em]
  \item \textbf{Relevance is single-judge} and declared by subtype rather than
        per item. The proposal calls for multi-judge agreement; that is
        outstanding.
  \item \textbf{The \code{filename} comparator is inflated.} The fetcher names
        files after their class (\code{scn\_beach\_020.jpg}), so path matching
        succeeds far more often than on a phone, where names look like
        \code{IMG\_20231205\_123456.jpg}. Treat 0.299 p@1 as a loose upper bound
        on the trivial floor, not the real floor.
  \item \textbf{\code{doctype} rules do not cover \code{bill}, \code{book},
        \code{handwritten} or \code{general\_doc}} --- 72 of 287 documents are
        mistagged. The rules were written when those classes had no data.
  \item \textbf{No screenshots}, though the proposal names them as a gap in
        existing benchmarks (20--40\% of a real gallery). Deferred.
  \item \textbf{Synthetic documents are rendered, not photographed.} The
        0.418 vs 0.592 CER gap bounds the optimism this introduces; the nine
        held-out real specimens are the check.
  \item \textbf{X-ray queries misroute} and handwriting retrieval fails
        outright --- both visible in \S\ref{sec:modules}, neither yet addressed.
\end{enumerate}

\end{document}
